\section{Gegnerwerte}

\subsection{Die Umgebung \texttt{Wertekasten}}

Die Umgebung \verb|Wertekasten| erzeugt einen standardisierten Werteblock
für Gegner.

Sie basiert auf der Umgebung \verb|FlexibleBox| und stellt innerhalb des
Blocks zusätzliche Befehle zur Verfügung:

\begin{itemize}
\item \verb|\wert|
\item \verb|\waffe|
\item \verb|\ruestung|
\end{itemize}

\subsubsection*{Syntax}

\begin{Verbatim}
\begin{Wertekasten}{Name}

...

\end{Wertekasten}
\end{Verbatim}

Der Parameter gibt den Titel der Box an.

\subsubsection*{Beispiel}

\paragraph{Quellcode}

\begin{Verbatim}
\begin{Wertekasten}{Goblin}

\wert{MU}{10}
\wert{KL}{8}
\wert{KK}{9}

\end{Wertekasten}
\end{Verbatim}

\paragraph{Ergebnis}

\begin{Wertekasten}{Goblin}

\wert{MU}{10}
\wert{KL}{8}
\wert{KK}{9}

\end{Wertekasten}


\subsection{Der Befehl \texttt{\textbackslash wert}}

Der Befehl \verb|\wert| erzeugt ein einzelnes Wertepaar.

Er ist nur innerhalb einer \verb|Wertekasten|-Umgebung verfügbar.

\subsubsection*{Syntax}

\begin{Verbatim}
\wert{Name}{Wert}
\end{Verbatim}

\subsubsection*{Parameter}

\begin{center}
\begin{tabular}{|l|l|}
\hline
Parameter & Bedeutung\\
\hline
Name & Bezeichnung des Wertes\\
\hline
Wert & eigentlicher Wert\\
\hline
\end{tabular}
\end{center}

\subsubsection*{Beispiel}

\paragraph{Quellcode}

\begin{Verbatim}
\begin{Wertekasten}{Ork}

\wert{MU}{12}
\wert{LE}{35}

\end{Wertekasten}
\end{Verbatim}

\paragraph{Ergebnis}

\begin{Wertekasten}{Ork}

\wert{MU}{12}
\wert{LE}{35}

\end{Wertekasten}


\subsection{Der Befehl \texttt{\textbackslash waffe}}

Der Befehl \verb|\waffe| erzeugt eine vollständige Waffenzeile.

Er trägt automatisch die Werte für

\begin{itemize}
\item AT (Attacke)
\item PA (Parade)
\item TP (Trefferpunkte)
\item INI (Initiative)
\end{itemize}

ein.

\subsubsection*{Syntax}

\begin{Verbatim}
\waffe
    {Name}
    {AT}
    {PA}
    {TP}
    {INI}
\end{Verbatim}

\subsubsection*{Beispiel}

\paragraph{Quellcode}

\begin{Verbatim}
\begin{Wertekasten}{Ork-Krieger}

\waffe
    {Axt}
    {12}
    {8}
    {1W6+3}
    {10}

\end{Wertekasten}
\end{Verbatim}

\paragraph{Ergebnis}

\begin{Wertekasten}{Ork-Krieger}

\waffe
    {Axt}
    {12}
    {8}
    {1W6+3}
    {10}

\end{Wertekasten}


\subsection{Der Befehl \texttt{\textbackslash ruestung}}

Der Befehl \verb|\ruestung| erzeugt eine Rüstungszeile.

Automatisch dargestellt werden:

\begin{itemize}
\item Rüstungsname
\item RS-Wert (Rüstungsschutz)
\end{itemize}

\subsubsection*{Syntax}

\begin{Verbatim}
\ruestung
    {Rüstung}
    {RS}
\end{Verbatim}

\subsubsection*{Beispiel}

\paragraph{Quellcode}

\begin{Verbatim}
\begin{Wertekasten}{Ork-Krieger}

\ruestung
    {Lederrüstung}
    {2}

\end{Wertekasten}
\end{Verbatim}

\paragraph{Ergebnis}

\begin{Wertekasten}{Ork-Krieger}

\ruestung
    {Lederrüstung}
    {2}

\end{Wertekasten}
